\section{\translate{Introduktion}}\label{sec:intro}
Projektet syftar till att undersöka hur organisatorisk information kan representeras för att med hjälp av Large Language Models (LLM:er) effektivisera intern informationssökning inom en organisation.

Projektet görs i samarbete med företaget Knightec Group AB.
\subsection{\translate{Bakgrund och problemmotivering}}\label{subsec:background}
\noindent 
I många organisationer är data lagrad och hanterad i flera separata system, 
vilket leder till att information isoleras i så kallade datasilos. 
Dessa kan bestå av affärssystem, kundhanteringssystem och andra interna databaser.

Denna spridning gör det svårt för anställda att få en samlad bild av verksamheten, då information manuellt behöver sökas upp och samlas in
från flera olika system\cite{2}.

Problemet är inte begränsat till enskilda organisationer, utan förekommer inom många sorters verksamheter.
Enligt en studie från International Data Corporation (IDC) citerad i \cite{armedia} visas det att en kunskapsarbetare 2013 spenderade i genomsnitt upp till 2,5 timmar per dag att söka information, vilket motsvarar ungefär 30\% av en vanlig arbetsdag.
Globalt uppskattas det leda till att företag förlorar stora summor pengar varje år på grund av denna ineffektivitet i informationshanteringen. Det leder inte bara till dessa ekonomiska förluster,
utan påverkar även de psykologiska aspekterna hos medarbetare, exempelvis frustration hos medarbetare när en bra idé dör på grund av den långa tid det tar att hitta stödjande data\cite{2}.

Artificiell intelligens har blivit en central del av många organisationers strategi för att hantera och utnyttja data.
Trots detta når många AI-initiativ aldrig bortom prototypstadiet. En avgörande orsak till detta är just dessa datasilos.
När information är uppdelad på detta sätt försvåras AI-systemens möjligheter att utnyttja informationen effektivt, eftersom AI-systemen är beroende av sammanhängande och kontextuell data\cite{3}.

Ett möjligt sätt att effektivisera informationssökningen och lösa problemet med svåråtkomlig information är genom grafbaserade kunskapssystem. I dessa kunskapssystem representeras informationen i form av noder och relationer i en grafstruktur.
Detta skapar en mer sammanhängande och kontextuell representation av data samlat på ett ställe antingen ovanpå organisationens befintliga system, eller genom att ersätta delar av dem.
Dessa kunskapsgrafer kan sedan användas av AI-system för att förstå relationer och navigera till relevant information på ett mer effektivt sätt\cite{1}.

\subsection{\translate{Övergripande syfte}}\label{subsec:aim}
\noindent

Syftet med arbetet är att undersöka hur organisatorisk information kan representeras för att effektivisera informationssökning med hjälp av LLM:er.
Detta görs genom att utveckla och jämföra hur LLM:en kan utnyttja en grafbaserad strukturering av data mot en lösning där LLM:en hämtar data direkt från de externa systemens API:er.

\subsection{\translate{Konkreta och verifierbara mål}}\label{subsec:researchquestion}
För att arbetet ska kunna utvärderas på ett verifierbart sätt har följande verifierbara mål satts upp.

\begin{enumerate}
  \item \textbf{Effektivitet i informationssökning } \\
   Hur lång tid tar det för en LLM att hitta och hämta relevant information i ett grafbaserat system jämfört med en API-baserad lösning?
  \item \textbf{Kvalitet på genererade svar} \\
   Hur ofta genereras relevanta och korrekta svar?
  \item \textbf{Resursanvändning vid informationshämtning} \\
   Hur många verktygsanrop behövs för att hitta informationen?
\end{enumerate}

\subsection{\translate{Avgränsningar}}\label{subsec:scope}
\noindent
För att arbetet ska vara genomförbart inom tidsramen har en del avgränsningar satts upp.

Arbetet inkluderar:
\begin{itemize}
  \item {Skapande av en kunskapsgraf som lagras i en grafdatabas, backend med Model Context Protocol (MCP) server för att låta LLM:en prata med grafdatabasen och användning av ett användargränssnitt med integrerad LLM för att tolka användarinmatning}
  \item {Skapande av ett system som använder den skapade MCP-servern för att prata direkt med olika systems API:er istället för kunskapsgrafen}
  \item {Tester och jämförelse av de punkter från \ref{subsec:researchquestion}}
\end{itemize}

Arbetet inkluderar inte:
\begin{itemize}
  \item{Skapande och utvärdering av andra typer av strukturer för att representera företagsinformation än kunskapsgrafer.}
  \item{Skapande av ett system för automatisk överföring av data från de externa systemen till grafdatabasen.}
\end{itemize}

\subsection{Användning av AI verktyg}
Artificiell intelligens används till att skapa och felsöka kod i arbetet, bidra med idéer och lösningar på problem som uppstår och som hjälpmedel för att formulera och förbättra text i rapporten.


\subsection{\translate{Översikt}}\label{subsec:outline}
\noindent

\texttt{Kapitel 2.} Teorin ger läsaren en förståelse för de olika bibliotek och verktyg som använts för att skapa Proof Of Concept (PoC) systemet.

\texttt{Kapitel 3.} Metoden tar upp hur genomförandet av arbetet har sett ut, det vill säga de arbetsmetoder som följts.

\texttt{Kapitel 4.} Förstudien redogör för de studier som genomförts för att identifiera problemen, lösningar och de tekniska verktyg som används.

\texttt{Kapitel 5.} Implementationskapitlet går igenom hur systemet har implementerats, fungerar och vilka tekniker som använts.

\texttt{Kapitel 6.} Resultat presenterar resultaten av implementationen och de verifierbara målen som satts upp.

\texttt{Kapitel 7.} I diskussionen diskuteras resultaten, metodvalen och begränsningarna i arbetet. Här diskuteras även möjliga förbättringar till framtida arbeten.

\texttt{Kapitel 8.} Slutsatsen ger en sammanfattning av arbetet och diskuterar uppnådda mål, samt ger förslag på framtida arbeten och utveckling av systemet.

